\documentclass[aspectratio=169]{beamer}

% Use Sorbonne theme
\usepackage{sorbonne_theme}

% Title page information
\title{Econometrics}
\subtitle{Binary Choice and Selection Models}
\author{Jose Angel Garcia Sanchez}
\institute{Université Paris 1 Panthéon-Sorbonne}
\date{} % left empty on purpose: beamer would otherwise print today's date
% Set the logo
\sorbonnelogo{imagen.png}

\begin{document}

% Title page (without logo in corner)
{
\setbeamertemplate{logo}{}
\begin{frame}
\titlepage
\end{frame}
}

% Table of contents
\begin{frame}{Outline}
\tableofcontents
\end{frame}

% ===============================
% SECTION 1: Maximum Likelihood
% ===============================
\section{Maximum Likelihood}

\begin{frame}{Presentation}
\begin{itemize}
    \item An ordinary linear model can be estimated via OLS
    \item But there exists another method: \textit{Maximum Likelihood}
    \item We write the \textit{likelihood} of the sample, which we maximize with respect to the parameters we'd like to estimate
    \item We then get the optimal parameters
    \item In the OLS method, we minimized the sum of squared residuals: it is again an optimization process
\end{itemize}
\end{frame}

\begin{frame}{The model (simple case)}
\begin{itemize}
    \item Let's consider an ordinary linear model $y_i = a + bx_i + u_i$
    \item Assume that the error term $u_i$ follows a normal distribution $N(0, \sigma^2)$
    \item We consider $a$, $b$, and $x$ as fixed
    \item Thus $y_i$ follows a normal distribution as well: $N(a + bx_i, \sigma^2)$
    \item The ``likelihood'' of the sample is the joint probability (density) of the observed sample, seen as a function of the parameters
\end{itemize}
\end{frame}

\begin{frame}{Estimation}
\begin{itemize}
    \item The distribution of $y$ depends on parameters $a$ and $b$ that are unknown
    \item If we think that the sample was correctly sampled, then parameters $a$ and $b$ are such that the probability of having sampled this particular data is maximum
    \item In other words, these parameters should maximize the likelihood of the sample
    \item This is the ML method: Maximum Likelihood
\end{itemize}
\end{frame}

\begin{frame}{Computation}
\begin{itemize}
    \item We assume that individuals in the sample are independent
    \item The likelihood of the sample is thus the product of the likelihood of each individual
    \item For this reason, we usually take its log to get a sum (easier to handle)
    \item We thus deal with the ``log-likelihood''
    \item Parameters $a$ and $b$ are obtained by maximizing the log-likelihood, which is the same as maximizing the likelihood since the log function is strictly increasing
    \item In practice: usual maximization techniques, or iterative methods
    \item With ML, we get the same estimates of $a$ and $b$ as with OLS, with $\hat{\sigma}^2$ biased downwards ($\hat{\sigma}^2_{ML} = SSR/N$)
\end{itemize}
\end{frame}

\begin{frame}{Application}
\begin{itemize}
    \item $\forall i$, $y_i \hookrightarrow N(a + bx_i, \sigma^2)$, with density $f(y_i) = \frac{1}{\sigma\sqrt{2\pi}} \exp\left(-\frac{1}{2\sigma^2}(y_i - a - bx_i)^2\right)$
    \item So the likelihood of the sample is $L = \prod f(y_i)$, and the log-likelihood is:
    \[LL = \log(L) = \sum \log(f(y_i)) = -\frac{N}{2}\log(2\pi) - N\log\sigma - \frac{1}{2\sigma^2}\sum(y_i - a - bx_i)^2\]
    \item In the end: $\max_{(a,b)} LL \Longleftrightarrow \min_{(a,b)} \sum(y_i - a - bx_i)^2$: in this special case, OLS and ML estimators of $a$ and $b$ are identical
    \item Similarly, we find $\hat{\sigma}^2_{ML} = \frac{SSR}{N}$, which is biased downwards (we found earlier that an unbiased estimator was $\hat{\sigma}^2_{OLS} = \frac{SSR}{N-k}$)
    \item However, although biased, it is consistent, since in a large sample $N$ and $N-k$ are almost the same
\end{itemize}
\end{frame}

\begin{frame}{Maximum Likelihood: General Properties}
Under regularity conditions:
\begin{itemize}
    \item ML estimators can be biased
    \item But they are consistent
    \item Asymptotically normally distributed
    \item Asymptotically efficient (they reach the Cramér-Rao bound: smallest asymptotic variance among consistent asymptotically normal estimators)
    \item Caution: the likelihood must be well specified
\end{itemize}
\end{frame}

\begin{frame}{Tests}
No OLS tests any more, since we are no longer in the OLS framework. Assume we want to test $q$ linear restrictions on parameter $\theta$: $R\theta = r$. The possibilities, which are asymptotically equivalent, are:

\begin{itemize}
    \item Wald test: estimate parameter $\theta$ by ML, then check whether $R\hat{\theta} - r$ is close to zero, using its asymptotic variance-covariance matrix (this procedure is close to the tests we are used to)
    \item Likelihood ratio test: estimate the model with and without the constraints, and check whether the difference between the 2 LL is significantly different from 0: $LR = 2(LL_{unconstrained} - LL_{constrained}) \hookrightarrow \chi^2_q$ under $H_0$
    \item Lagrange multiplier test: estimate the model with the restrictions, then check whether the first order condition of the unrestricted model is significantly violated
\end{itemize}
\end{frame}

% ===============================
% SECTION 2: Binary Choice Models
% ===============================
\section{Models of Discrete Choice}

\begin{frame}{Binary Choice}
\begin{itemize}
    \item Example: for individual $i$, $y_i = 1$ if unemployed, $y_i = 0$ if employed
    \item We want to explain $y$ by age, education, etc.
    \item OLS fits a line, but here there is no real scatterplot since $y$ can take only 2 values, 0 and 1
    \item Plus, many other problems, the main one being that predictions will be out of range
    \item So the preferred method is ML
\end{itemize}
\end{frame}

\begin{frame}{Discrete Choice Models and OLS (1)}
\begin{itemize}
    \item Example: for individual $i$, $y_i = 1$ if unemployed, $y_i = 0$ if employed
    \item We want to explain $y$ by age and education (the \textit{linear probability model}):
    \item $y_i = \alpha + \beta\,\text{age}_i + \gamma\,\text{education}_i + \varepsilon_i = x_i b + \varepsilon_i$
    \item So $E(y_i|x_i) = x_i b$
    \item And we know that $E(y_i|x_i) = 1 \times P(y_i = 1|x_i) + 0 \times P(y_i = 0|x_i)$, so $P(y_i = 1|x_i) = x_i b$
    \item $x_i b$ should lie between 0 and 1, which imposes strong restrictions on $x$ and $b \Rightarrow$ inconvenient
\end{itemize}
\end{frame}

\begin{frame}{Discrete Choice Models and OLS (2)}
\begin{itemize}
    \item Furthermore, $\varepsilon_i$ can take only 2 values: $-x_i b$ with probability $1 - x_i b$, or $1 - x_i b$ with probability $x_i b$
    \item So that: $V(\varepsilon_i) = x_i b(1 - x_i b)$
    \item There is heteroskedasticity
\end{itemize}
\end{frame}

\begin{frame}{Discrete Choice Models and OLS (3)}
\begin{itemize}
    \item Heteroskedasticity could be dealt with using White's standard errors (see lecture on FGLS)
    \item Non-normality could be dealt with using a large sample size (Central Limit Theorem)
    \item But the main problem is that we need $E(y_i|x_i)$ to belong to $[0,1]$
    \item For these various reasons, OLS is inappropriate for predicting probabilities (although the linear probability model with robust standard errors is still often used to estimate average effects)
\end{itemize}
\end{frame}

\begin{frame}{Solution}
\begin{itemize}
    \item We act as if there was an unobserved continuous variable $y^*$ such that $\forall i$, $y_i^* > 0 \Rightarrow y_i = 1$ and $y_i^* \leq 0 \Rightarrow y_i = 0$
    \item This new variable $y^*$ is continuous: we can thus refer to our usual model, with some adaptations
\end{itemize}
\end{frame}

\begin{frame}{The Model}
\begin{itemize}
    \item Assume that $y^* = X\beta + u$
    \item Let $F$ be the cumulative distribution function (cdf) of $u$, with a distribution symmetric around 0
    \item $P(y_i = 1) = P(y_i^* > 0) = P(X_i\beta + u_i > 0) = P(u_i > -X_i\beta) = P(u_i < X_i\beta) = F(X_i\beta)$
    \item Usual functions for $F$ are the normal (\textit{Probit model}) or the logistic (\textit{Logit model})
\end{itemize}
\end{frame}

\begin{frame}{Distributions}
\begin{itemize}
    \item Normal: $F(x) = \Phi(x) = \int_{-\infty}^{x} \frac{1}{\sqrt{2\pi}}e^{-t^2/2}dt$
    \item Logistic: $F(x) = \frac{1}{1+e^{-x}} = \frac{e^x}{1+e^x}$
    \item These 2 are very close, and give similar results
    \item Expectation: 0 for both distributions
    \item Variance: 1 for the standard normal, $\pi^2/3$ for the logistic
    \item Characteristics of the \textit{Logit}: extreme events have higher probability (heavier tails); parameters are interpreted more easily (odds-ratios, see below)
\end{itemize}
\end{frame}

\begin{frame}{Log-likelihood (1)}
Likelihood function for the entire sample:

\[L = \prod P[y_i = 1|x_i, \beta]^{y_i} P[y_i = 0|x_i, \beta]^{1-y_i}\]

So that when $y_i = 0$, only the right-hand-side part is there, and when $y_i = 1$, only the left-hand-side part is there.

And we have $P[y_i = 1|x_i, \beta] = P[y_i^* > 0|x_i, \beta] = F(x_i\beta)$, so

\[LL = \sum y_i \log(F(x_i\beta)) + \sum(1 - y_i)\log(1 - F(x_i\beta))\]

Maximizing this log-likelihood amounts to differentiating this expression with respect to parameter $\beta$ and setting it to zero (it can be shown that it is globally concave for the Probit and the Logit).
\end{frame}

\begin{frame}{Log-likelihood (2)}
\[LL = \sum y_i \log(F(x_i\beta)) + \sum(1 - y_i)\log(1 - F(x_i\beta))\]

The first order condition is the following (derivative wrt $\beta$ is 0): $\frac{\partial LL}{\partial \beta} = 0$, which in fact is a column vector of derivatives (derivative of $LL$ wrt $\beta_1$, $\beta_2$, etc.)

\[\frac{\partial LL}{\partial \beta} = \sum\left[\frac{y_i - F(x_i\beta)}{F(x_i\beta)(1 - F(x_i\beta))}f(x_i\beta)\right]x_i' = 0\]

For the logistic distribution, $F(x) = \frac{e^x}{1+e^x}$ and $f(x) = \frac{e^x}{(1+e^x)^2} = F(x)(1-F(x))$, so:

\[\frac{\partial LL}{\partial \beta} = \sum\left(y_i - \frac{\exp(x_i\beta)}{1 + \exp(x_i\beta)}\right)x_i' = 0\]

The optimal value for $\beta$ is found with optimization techniques (iterative).
\end{frame}

\begin{frame}{Remark: Scaling}
\begin{itemize}
    \item Recall that $P(y_i = 1) = P(u_i < X_i\beta) = F(X_i\beta)$: this holds only if $u$ has a standardized distribution (of variance 1 for the Probit)
    \item $u$ could in fact have \textit{any} standard deviation $\sigma$; in that case: $P(u_i < X_i\beta) = P(u_i/\sigma < X_i\beta/\sigma) = F(X_i\beta/\sigma)$
    \item So the estimate we compute is really $\beta/\sigma$: there is no way to identify the two separately
    \item The convention is to set $\sigma = 1$ for the Probit model, and $\sigma = \pi/\sqrt{3}$ (the standard deviation of the logistic distribution) for the Logit model
    \item This multiplicative scaling changes neither the signs nor the ratios $\beta_j/\beta_k$ of parameters, nor the predictions
    \item The only consequence is that Probit and Logit parameters are scaled differently: $\hat{b}_{logit} \simeq 1.6\,\hat{b}_{probit}$ in practice ($\pi/\sqrt{3} \simeq 1.8$ if we only match the variances)
\end{itemize}
\end{frame}

\begin{frame}{Remark: Predicted and Actual Frequency}
\begin{itemize}
    \item $\hat{p}_i = \hat{P}(y_i = 1) = F(x_i\hat{\beta})$
    \item For the Logit, $\hat{p}_i = 1/(1 + \exp(-x_i\hat{\beta})) = \exp(x_i\hat{\beta})/(1 + \exp(x_i\hat{\beta}))$
    \item $\hat{\beta}$ is the solution of the first order condition, so that we get:
\end{itemize}

\[\sum\left(y_i - \frac{\exp(x_i\hat{\beta})}{1 + \exp(x_i\hat{\beta})}\right)x_i' = 0\]

\[\sum(y_i - \hat{p}_i)x_i' = 0 \quad\Longleftrightarrow\quad \sum \hat{p}_i x_i' = \sum y_i x_i'\]

If there is a constant term in $x$, then $\sum \hat{p}_i = \sum y_i$: the predicted frequency is exactly equal to the actual frequency. This holds approximately for the Probit model.
\end{frame}

\begin{frame}{Interpreting Parameters (1)}
Let $x_i$ be a row vector of observations (corresponding to individual $i$) and $\beta$ be a column vector of parameters.

\begin{itemize}
    \item $P(y_i = 1) = F(x_i\beta)$
    \item Parameters enter non-linearly in the expression
    \item If we differentiate $F(x_i\beta)$ with respect to variable $x_k$, calling $f$ the derivative of $F$, we get:
\end{itemize}

\[\frac{\partial F(x_i\beta)}{\partial x_{i,k}} = f(x_i\beta)\beta_k\]

With $f$ either the standard normal density $\phi$ (\textit{probit}) or the logistic density $\exp(x)/(1+\exp(x))^2$ (\textit{logit}).

The effect of a change in $x_{i,k}$ depends on the values of $x_i$, so it is different for each person.
\end{frame}

\begin{frame}{Interpreting Parameters (2)}
\[\frac{\partial F(x_i\beta)}{\partial x_{i,k}} = f(x_i\beta)\beta_k\]

\begin{itemize}
    \item The \textit{sign} of the effect of a change in $x_{i,k}$ can however be determined
    \item Since $f > 0$, it is the sign of $\beta_k$
\end{itemize}
\end{frame}

\begin{frame}{Marginal Effects}
\begin{itemize}
    \item We still miss a way to measure the effect of $X$ on $y$: the \textit{marginal effects}
    \item Consider the average person in the sample and compute the prediction $\hat{y}_0 = F(\bar{X}\hat{\beta})$
    \item Then increase variable $x_1$ by one unit, and compute the prediction again: $\hat{y}_1$
    \item The marginal effect of $x_1$ (at the mean) is $\hat{y}_1 - \hat{y}_0$: by how much the probability of $[y = 1]$ increases when $x_1$ increases by one unit (e.g.\ one more year of education, switching from male to female if gender is a dummy)
    \item Always mention that this was computed for the average person; it can also be computed for specific individuals (e.g.\ the average 40-year-old white-collar worker)
    \item Alternative: the \textit{average marginal effect} (AME), i.e.\ the individual effects averaged over the whole sample
    \item statsmodels: \texttt{res.get\_margeff(at="mean")}, or \texttt{at="overall"} for the AME
\end{itemize}
\end{frame}

\begin{frame}{Logit vs. Probit}
\begin{itemize}
    \item Since the two distributions are quite similar, predictions will be similar as well
    \item However, values of parameters differ: this is because of the different scaling
    \item We usually get $\hat{b}_{logit} \simeq 1.6\, \hat{b}_{probit}$
    \item Matching only the variances would give: $\hat{b}_{logit} \simeq (\pi/\sqrt{3}) \hat{b}_{probit} \simeq 1.8\, \hat{b}_{probit}$
    \item Again, this does not mean that the impact of explanatory variables is higher with the Logit model
\end{itemize}
\end{frame}

\begin{frame}{Logit: Odds-ratios (1)}
\begin{itemize}
    \item Parameters $\beta$ do not have a linear impact: we can compute marginal effects, but the Logit provides more information
    \item Recall that $P(y_i = 1) = P(y_i^* > 0) = F(X_i\beta) = \frac{1}{1+e^{-X_i\beta}}$
    \item I.e.\ $\forall i$, $P(y_i = 1) = p_i$ and $P(y_i = 0) = 1 - p_i$
    \item Let's define the \textit{odds} $c_i$ as $c_i = \frac{p_i}{1-p_i}$
    \item $[y_i = 1]$ is $c_i$ times more likely to happen than $[y_i = 0]$
    \item We have $p_i = \frac{1}{1+e^{-X_i\beta}} = \frac{e^{X_i\beta}}{1+e^{X_i\beta}}$
    \item And $1 - p_i = \frac{1}{1+e^{X_i\beta}}$
    \item So $c_i = e^{X_i\beta}$
\end{itemize}
\end{frame}

\begin{frame}{Logit: Odds-ratios (2)}
\begin{itemize}
    \item We thus get $c_i = e^{X_i\beta}$
    \item But $X_i\beta = \beta_0 + \sum_j x_{i,j}\beta_j$
    \item So $c_i = e^{\beta_0 + \sum_j x_{i,j}\beta_j} = e^{\beta_0}\prod_j e^{x_{i,j}\beta_j}$
    \item If variable $x_{i,k}$ increases by 1 unit, $c_i$ is multiplied by $e^{\beta_k}$
    \item If variable $x_k$ increases by 1 unit, it multiplies the \textit{odds} that $[y_i = 1]$ happens by $e^{\beta_k}$
    \item We call \textit{odds-ratio} associated with variable $x_k$ the value $e^{\beta_k}$ (in Python: \texttt{np.exp(res.params)})
\end{itemize}
\end{frame}

\begin{frame}{Goodness of Fit (1)}
\begin{itemize}
    \item Pseudo-$R^2$ (McFadden): $R^2 = \frac{LL_{fit} - LL_0}{LL_{max} - LL_0} = 1 - \frac{LL_{fit}}{LL_0}$ (since $LL_{max} = 0$)
    \item $LL_{fit}$: log-likelihood of the model
    \item $LL_0$: log-likelihood of a model with only a constant, where the probability is estimated by the proportion of ones in the sample
    \item $LL_{max}$: maximum log-likelihood attainable. In our case, 0, because a model that predicts the observed values perfectly has likelihood 1, and $\log(1) = 0$
    \item statsmodels reports it as ``Pseudo R-squ.''
\end{itemize}
\end{frame}

\begin{frame}{Goodness of Fit (2)}
\begin{itemize}
    \item We can also evaluate the goodness of fit of the model by comparing correct and incorrect predictions
    \item Predicted outcomes: since $0 \leq \hat{p}_i \leq 1$, we need to choose a cut-off point $c$ that will determine if the prediction will be 0 or 1 (usually: $c = 1/2$)
    \item \textit{Sensitivity} (sensibilité): true positive rate in predictions
    \item \textit{Specificity} (spécificité): true negative rate in predictions
    \item ROC (Receiver Operating Characteristic) curve: plots the fraction of $y = 1$ values correctly classified (true positives) against the fraction of $y = 0$ values incorrectly classified (false positives) as $c$ varies, i.e.\ \textit{sensitivity} against \textit{(1 $-$ specificity)}
\end{itemize}
\end{frame}

\begin{frame}{The ROC Curve}
\begin{itemize}
    \item Comes from signal detection theory, developed during WW2
    \item The accuracy of the classification rule is measured with the area under the ROC curve (AUC)
    \item The area represents the ability of the procedure to correctly classify individuals, and lies between 0.5 (random classification) and 1 (perfect classification)
    \item A random classification rule yields a diagonal ROC curve
    \item Trade-off between sensitivity and specificity: the more bowed the curve is, the better the model is
    \item The closer the ROC curve to the upper-left corner, the higher the predictive power
    \item The optimal cut-off point is determined by the analyst depending on the issue studied
\end{itemize}
\end{frame}

\begin{frame}{Could Heteroskedasticity be an Issue?}
\begin{itemize}
    \item In Probit and Logit models, homoskedasticity of $u$ is assumed in the latent variable framework (normalization: $\sigma = 1$)
    \item Recall that in the Probit model, only $\beta/\sigma$ can be identified
    \item Heteroskedasticity means that $\sigma$ is a function of some variables, say $\sigma_i = \exp(Z_i\gamma)$
    \item Remark 1: it is best not to have a constant in variables $Z$, because then testing $\gamma = 0$ amounts to testing $\sigma = \exp(0) = 1$, which means homoskedasticity
    \item Remark 2: variables $Z$ should be exogenous variables not already in the model (see Cameron \& Trivedi)
    \item In this case, we model $P(y_i = 1) = P(y_i^* > 0) = F(X_i\beta/\sigma_i)$ and plug in $\sigma_i = \exp(Z_i\gamma)$ in the likelihood
    \item Stata's \textit{hetprobit} provides this estimation, with a test of $\gamma = 0$; in Python, the likelihood can be written with statsmodels' \texttt{GenericLikelihoodModel}
\end{itemize}
\end{frame}

\begin{frame}{Could Endogeneity be an Issue?}
\begin{itemize}
    \item Just like OLS, Probit and Logit estimates are inconsistent if one or several explanatory variables are endogenous
    \item The usual latent variable model is written: $y^* = X\beta + x\gamma + u$, with $x$ endogenous
    \item In that case, given we find enough instruments $Z$, the endogenous variable should be instrumented: $x = Zb + Xc + v$
    \item To estimate this now 2-equation model in the maximum likelihood framework, we only need to postulate that $(u, v)$ follow a bivariate normal distribution and plug this into the likelihood
    \item Testing endogeneity of $x$ amounts to testing whether $u$ and $v$ are uncorrelated
    \item All this requires both normality and homoskedasticity of $u$ and $v$, so Probit is required (no Logit here)
\end{itemize}
\end{frame}

% ===============================
% SECTION 3: Sample Selection Models
% ===============================
\section{Sample Selection Models}

\begin{frame}{About Selection}
\begin{itemize}
    \item E.g.\ consider the interpretation of average scores over time of an achievement test
    \item A decline over time may be due to real deterioration in student knowledge, or it may just reflect a selection effect, i.e.\ more students have been taking the test over time and the new test takers are the relatively weaker students
    \item Selection can arise from self-selection (individuals choose to participate or not in a particular activity) or sample selection (the sample is not random: e.g.\ those who participate in the activity are oversampled)
    \item These problems are treated alike in Econometrics through \textit{sample selection models}
\end{itemize}
\end{frame}

\begin{frame}{The Heckit Model (or Generalized Tobit)}
\begin{itemize}
    \item Say a variable $y_2^*$ is partially observed ($y_2^*$ is the latent variable)
    \item Say we have another latent variable $y_1^*$, such that $y_2^*$ is observed only if $y_1^* > 0$
    \item E.g.\ health expenditures or wages
    \item $y_1$ is participation in the labor market: $y_1 = 0$ (does not participate) or $y_1 = 1$ (participates)
    \item $y_2$ is the observed wage, which cannot be observed if the person does not participate in the labor market: that person may have a positive potential wage $y_2^*$, but we only observe $y_1 = 0$
    \item We'll see here why the Probit model is useful for the estimation of the 1st step equation
    \item Note: in the \textit{standard} Tobit model, $y_2$ is observed only if it is above (or below) a certain threshold
\end{itemize}
\end{frame}

\begin{frame}{The Model}
Latent variables: $y_1^* = x_1b_1 + u_1$ and $y_2^* = x_2b_2 + u_2$

\begin{itemize}
    \item Participation equation: $y_1 = 1$ if $y_1^* > 0$, $y_1 = 0$ otherwise
    \item Outcome equation: $y_2 = y_2^*$ if $y_1 = 1$, $y_2$ is missing otherwise
\end{itemize}

\[\begin{pmatrix} u_1 \\ u_2 \end{pmatrix} \hookrightarrow N\left(\begin{pmatrix} 0 \\ 0 \end{pmatrix}; \begin{pmatrix} 1 & \sigma_{1,2} \\ \sigma_{1,2} & \sigma_2^2 \end{pmatrix}\right)\]

The Logit model could not be used here for the 1st equation, since the error terms are jointly normal: we have to consider the Probit model.
\end{frame}

\begin{frame}{The Likelihood}
\[L = \prod[P(y_{1,i}^* \leq 0)]^{1-y_{1,i}}[f(y_{2,i}|y_{1,i}^* > 0) \times P(y_{1,i}^* > 0)]^{y_{1,i}}\]

\begin{itemize}
    \item First term: discrete contribution when $y_{1,i}^* \leq 0$
    \item Second term: continuous contribution when $y_{1,i}^* > 0$
\end{itemize}
\end{frame}

\begin{frame}{Conditional Means}
\[E(y_2|x, y_1^* > 0) = E(x_2 b_2 + u_2|x_1 b_1 + u_1 > 0) = x_2 b_2 + E(u_2|u_1 > -x_1 b_1)\]

In the normal case, it can be shown that:

\[E(y_2|x, y_1^* > 0) = x_2 b_2 + \sigma_{1,2}\lambda(x_1 b_1)\]

With $\lambda$ the \textit{inverse Mills ratio}: $\lambda(z) = \frac{\phi(z)}{\Phi(z)}$

\begin{itemize}
    \item If $u_1$ and $u_2$ are independent, then $E(u_2|u_1 > -x_1 b_1) = 0$ and OLS are consistent in the second equation; otherwise, they are not: we have to take into account the sample selection bias $\sigma_{1,2}\lambda(x_1 b_1)$
    \item Error terms represent unobservable heterogeneity: it is likely that individual heterogeneity has an impact both on the decision to participate and on the outcome (e.g.\ health care expenditures)
    \item The sign of $\operatorname{corr}(u_1, u_2)$ is the sign of $\sigma_{1,2}$
\end{itemize}
\end{frame}

\begin{frame}{The Heckman Two-Step Estimator}
\[E(y_2|x, y_1^* > 0) = x_2 b_2 + \sigma_{1,2}\lambda(x_1 b_1)\]

Idea: first estimate $\lambda$, then estimate the model replacing $\lambda$ by its estimate $\hat{\lambda}$

\begin{enumerate}
    \item Probit regression of $y_1$ on $x_1$ (indeed, $P(y_1^* > 0) = \Phi(x_1 b_1)$)
    \item Compute $\hat{\lambda} = \frac{\phi(x_1 \hat{b}_1)}{\Phi(x_1 \hat{b}_1)}$
    \item OLS regression of $y_2$ on $x_2$ and $\hat{\lambda}$, on the selected sample only ($y_1 = 1$)
    \item The coefficient of $\hat{\lambda}$ estimates $\sigma_{1,2}$: its sign gives the sign of the correlation between the 2 error terms, and if it is not significant, the 2 equations can be considered independent
\end{enumerate}
\end{frame}

\begin{frame}{Remarks on the two-step estimator}
\begin{itemize}
    \item Two-step Heckman estimates are consistent
    \item This method is easy and fast to implement
    \item But the usual OLS standard errors of the second step are wrong ($\hat{\lambda}$ is itself estimated): they must be corrected
    \item There is also an efficiency loss (not appropriate in small samples)
    \item A more widely used method is Maximum Likelihood on the full model
    \item Notice that the use of a Logit would not be appropriate here
\end{itemize}
\end{frame}

\begin{frame}{Identification Considerations}
\begin{itemize}
    \item Theoretically, the exact same regressors can appear in both equations, because $\lambda(\cdot)$ is non-linear: no strict multicollinearity
    \item However, in that case, the model is close to unidentified, because the inverse Mills ratio is almost linear over a wide range of values
    \item This leads to a great instability of parameters (see slides on near-perfect multicollinearity)
    \item Thus, at least one exclusion restriction is usually required: at least one variable in the first equation (selection) should be absent from the second equation (conditional outcome)
\end{itemize}
\end{frame}

\begin{frame}{Selection on observables vs.\ unobservables}
\begin{itemize}
    \item If we find out that there is in fact no correlation between the 2 equations (check $\rho = \operatorname{corr}(u_1, u_2) = \sigma_{1,2}/\sigma_2$), they can be estimated independently through a Two-Part model (see next slide)
    \item If there is no correlation: selection is said to be made on observables, because it is fully explained by regressors
    \item If there is correlation: selection is said to be made on unobservables, because it is partly explained by error terms, which comprise all the unobservable information
    \item The Two-Part model can also be used if the correlation of the inverse Mills ratio with regressors is too strong (near-perfect multicollinearity)
\end{itemize}
\end{frame}

\begin{frame}{The Two-Part Model}
This model intentionally does not take into account potential correlation between the error terms of the 2 equations

\begin{enumerate}
    \item Participation equation: Probit or Logit that gives the probability to participate
    \item Conditional outcome equation: anything we like (linear model, count data, etc.) that gives the level of activity, given that the individual participates (warning: count models then need a distribution truncated at 0, such as the zero-truncated Poisson or negative binomial); so it is more flexible than the Generalized Tobit framework
\end{enumerate}

And the prediction of the model is the product of the two: $E(y) = P(y > 0) \times E(y|y > 0)$. Notice that we lose the assessment of the selection effect but gain in flexibility (no need to stick to Normal models).
\end{frame}

\end{document}
